% META: {"id": "1NSI-PM-RE-C2", "chapitre": "1NSI-PROJET-METHODES", "type_objet": "remediation", "capacites": ["BO-PREAMBULE-DEMARCHE-DE-PROJET"], "capacites_codes": ["C2"], "mode_creation": "ex_nihilo", "sources_inspiration": [], "status": "needs_review"}

%---------------------------------------------------------------
% FICHE DE REMEDIATION — C2 : un jalon se verifie, il ne se proclame pas
%---------------------------------------------------------------

\textbf{Ce qui bloque.} Tes jalons sont des intentions (« tout coder », « que ça marche »)
et non des étapes vérifiables : l'avancement reste à $0\,\%$ pendant des semaines, puis
saute à $100\,\%$ la veille, sans que personne sache où en est le projet. Et tu attends du
découpage qu'il garantisse l'absence de bug, ce qu'aucun jalon ne peut faire.

\textbf{À retenir.} Un jalon porte un critère de validation exécutable : un test qui passe,
une démonstration précise. Le pourcentage compte des jalons vérifiés ; il ne mesure ni la
charge de travail restante ni la qualité du code.

\begin{exercice}{1NSI-PM-RE-C2-EX1}{1}{12}
Une équipe a décrit son projet de gestionnaire de tâches ainsi :
\begin{python}
jalons = [
    {"nom": "Tout coder", "termine": False},
    {"nom": "Que ca marche", "termine": False},
]
\end{python}
Après trois semaines et $300$ lignes de code, \lstinline{avancement(jalons)} vaut toujours
$0{,}0$.
\begin{enumerate}
  \item Pourquoi ce pourcentage ne dit-il rien du travail accompli ? Que faudrait-il pour
        qu'un jalon puisse être déclaré terminé ?
  \item Réécrire ce projet en quatre jalons ordonnés, chacun avec un critère de validation
        que l'on peut exécuter ou montrer. Le premier doit être le cahier des charges, le
        dernier la présentation.
  \item Deux de ces quatre jalons sont terminés. Que valent \lstinline{avancement} et
        \lstinline{prochain_jalon} ? Ce $50\,\%$ garantit-il que la moitié du travail est
        faite, ou que le code déjà écrit est sans bug ?
\end{enumerate}
\end{exercice}

% BEGIN-VERIFY
% def avancement(jalons):
%     if len(jalons) == 0:
%         raise ValueError("au moins un jalon est requis")
%     nb_termines = sum(1 for j in jalons if j["termine"])
%     return nb_termines / len(jalons) * 100
% def prochain_jalon(jalons):
%     for j in jalons:
%         if not j["termine"]:
%             return j["nom"]
%     return None
% vagues = [{"nom": "Tout coder", "termine": False}, {"nom": "Que ca marche", "termine": False}]
% assert avancement(vagues) == 0.0
% assert prochain_jalon(vagues) == "Tout coder"
% jalons = [
%     {"nom": "Cahier des charges ecrit", "critere": "liste des fonctionnalites relue", "termine": True},
%     {"nom": "Ajout et affichage des taches", "critere": "ajouter 3 taches, les afficher", "termine": True},
%     {"nom": "Suppression et sauvegarde", "critere": "supprimer une tache, relancer, la retrouver absente", "termine": False},
%     {"nom": "Presentation repetee", "critere": "demonstration de 2 minutes", "termine": False},
% ]
% assert avancement(jalons) == 50.0
% assert prochain_jalon(jalons) == "Suppression et sauvegarde"
% END-VERIFY
